\documentclass[11pt,a4paper]{article}
\usepackage[utf8]{inputenc}
\usepackage[T1]{fontenc}
\usepackage[brazilian]{babel}
\usepackage[margin=2cm]{geometry}
\usepackage{lmodern}
\usepackage{amssymb}
\usepackage{xcolor}
\usepackage{array}
\usepackage{booktabs}
\usepackage{longtable}
\usepackage{enumitem}
\usepackage{titlesec}
\usepackage{fancyhdr}
\usepackage{parskip}
\usepackage[hidelinks]{hyperref}

\definecolor{brandgreen}{HTML}{2E5E3E}
\definecolor{lightbg}{HTML}{F5F2EC}

\titleformat{\section}{\large\bfseries\color{brandgreen}}{}{0pt}{}[{\color{brandgreen}\titlerule}]
\titleformat{\subsection}{\bfseries\color{brandgreen}}{}{0pt}{}

\pagestyle{fancy}
\fancyhf{}
\fancyhead[L]{\small Querência ERP --- Avaliação de usabilidade}
\fancyhead[R]{\small Roteiro do Condutor}
\fancyfoot[C]{\small\thepage}
\renewcommand{\headrulewidth}{0.4pt}

\newcommand{\caixa}[1]{%
  \noindent\colorbox{lightbg}{%
    \begin{minipage}{0.95\linewidth}
    \vspace{1mm}
    #1
    \vspace{1mm}
    \end{minipage}%
  }\vspace{3mm}
}

\begin{document}

\begin{center}
  {\LARGE\bfseries\color{brandgreen} Roteiro do Condutor}\\[2mm]
  {\large Sessões de avaliação de usabilidade --- Querência ERP}
\end{center}

\vspace{2mm}

\caixa{Este documento é o guia de ponta a ponta do moderador e \textbf{não é entregue ao
participante}. Os passos que o participante segue estão em 6 folhas separadas, um por fluxo
$\times$ papel: \textbf{fluxo-1-diretoria.pdf}, \textbf{fluxo-1-associado.pdf},
\textbf{fluxo-2-diretoria.pdf}, \textbf{fluxo-2-associado.pdf}, \textbf{fluxo-3-diretoria.pdf} e
\textbf{fluxo-3-associado.pdf} --- entregue ao participante só as 3 do seu papel, em sequência (F1,
F2, F3). Essas folhas não trazem falas do moderador, estado de kit nem comandos de verificação:
tudo isso está na Seção 7 deste roteiro.}

\section*{1. Objetivo da sessão}

Validar, com usuários reais das duas entidades parceiras, se o ecossistema digital qualifica os
processos administrativos identificados no diagnóstico --- em especial, se o fluxo mediado pela
entidade (RNF01) funciona como alternativa viável para participantes com menor familiaridade
digital. A sessão \textbf{não} busca validar a arquitetura técnica (já coberta pelos testes
automatizados do backend) --- busca validar a experiência de uso.

\section*{2. Formato: quatro blocos sequenciais (\textasciitilde35--45 min por participante)}

\begin{enumerate}[leftmargin=6mm]
  \item \textbf{Demonstração geral} (5 a 10 min) --- apresentar o sistema (painel e/ou
    app, conforme o papel do participante) de forma guiada, sem o participante tocar ainda,
    cobrindo navegação geral e terminologia da interface. Objetivo: nivelar familiaridade antes de
    medir desempenho em tarefa.
  \item \textbf{Execução das tarefas} (\textasciitilde15--20 min) --- o participante executa, sem
    ajuda do moderador a menos que trave por mais de \textasciitilde2 min, as tarefas do(s)
    fluxo(s) atribuídos ao seu papel. Cronometrar e anotar pontos de hesitação/erro, sem intervir
    preventivamente.
  \item \textbf{Questionário SUS} (\textasciitilde5 min) --- aplicado imediatamente após as
    tarefas, enquanto a experiência ainda está fresca.
  \item \textbf{Entrevista semiestruturada} (\textasciitilde10 min) --- roteiro de perguntas
    abertas (Seção 8.2 abaixo).
\end{enumerate}

\section*{3. Participantes}

2 a 4 participantes por entidade (4 a 8 no total), distribuídos em dois papéis:

\begin{itemize}[leftmargin=6mm]
  \item \textbf{Diretoria / responsável administrativo} --- testa os passos mediados pela
    entidade, no painel web.
  \item \textbf{Associado} --- testa os passos de autoatendimento, no aplicativo mobile.
\end{itemize}

Distribuição alvo por entidade: \textbf{1 a 2 da Diretoria + 1 a 2 Associados} (mínimo viável:
1 + 1). Com 4 pessoas, o ideal é 2 + 2 e pelo menos uma pessoa idosa ou de menor familiaridade
digital entre os Associados --- esse é o ponto que valida (ou não) RNF01 na prática.

\section*{4. Checklist antes de começar}

\begin{itemize}[leftmargin=6mm]
  \item[$\square$] Celular do moderador com o APK já instalado, pronto para ser entregue ao
    participante durante os cenários do papel Associado. Para o papel Diretoria, acesso ao painel
    web logado como administrador.
  \item[$\square$] Conectividade confirmada no local (Wi-Fi da sede ou dados móveis) --- o sistema
    depende de um backend real na nuvem, não funciona offline.
  \item[$\square$] \texttt{backend/scripts/verificar-tarefas.js} acessível (terminal com
    \texttt{cd backend} já configurado, \texttt{.env} com acesso ao banco).
  \item[$\square$] \textbf{Kit do participante conferido} (Seção 6 abaixo): personas, CPF
    pré-cadastrado, conta de teste com mensalidade em aberto, 2 ingressos pré-emitidos e associado
    inadimplente reservado --- tudo conferido \textbf{na véspera}, no painel, não no dia.
  \item[$\square$] Cronômetro e planilha/papel para registrar tempo, erros e intervenções.
  \item[$\square$] Formulário SUS pronto (digital ou impresso), conforme o papel do participante.
  \item[$\square$] Fichas de persona impressas (\texttt{ficha-participante.pdf}), reservadas em
    sequência (P01, P02\ldots).
\end{itemize}

\section*{5. Princípio de independência entre tarefas}

Com 2--4 participantes por entidade, \textbf{nenhuma tarefa pode depender do resultado de outra
tarefa nem da ação de outro participante}. Se um passo falhar (erro do participante, bug ou
rede), só aquela tarefa é perdida --- registra-se como \textit{não concluída} e a sessão segue.
Todo o estado de que uma tarefa precisa é preparado pelo moderador \textbf{antes} da sessão (o kit,
Seção 6), e não produzido por outra tarefa:

\begin{center}
\small
\begin{longtable}{p{4cm}p{5.3cm}p{4.7cm}}
\toprule
\textbf{Tarefa} & \textbf{Estado de que precisa} & \textbf{De onde vem} \\
\midrule
\endhead
F1 Diretoria (cadastrar) & Ficha de persona não usada & Ficha impressa (Pxx) \\
F1 Associado A (auto-cadastro) & Ficha de persona não usada & Ficha impressa (Pxx) \\
F1 Associado B (vincular conta) & CPF de associado ativo e sem conta de login & Pré-cadastrado
  pelo moderador no painel, antes do dia \\
F2 Associado (comprar ingresso) & Conta de teste ativa + evento publicado & Conta de teste do kit
  (login feito pelo moderador) \\
F2 Diretoria (validar entrada) & 2 ingressos já emitidos, em nome de associados fictícios &
  Emitidos pelo moderador antes; QR mostrado no celular do moderador \\
F3 Associado (pagar) & Conta de teste com mensalidade em aberto & Conta de teste do kit \\
F3 Diretoria (pagamento presencial) & 1 associado ativo e inadimplente & Reservado pelo moderador
  antes, anotado no kit \\
\bottomrule
\end{longtable}
\end{center}

Cada kit é de \textbf{uso único} (ingresso validado, mensalidade paga e persona cadastrada não
voltam ao estado inicial). Preparar 1 kit reserva por entidade além dos participantes confirmados.

\section*{6. Kits do participante (preencher na véspera)}

\subsection*{Kit --- papel Diretoria}

\begin{center}
\small
\begin{longtable}{p{7.5cm}p{3.3cm}p{2cm}}
\toprule
\textbf{Item} & \textbf{Preencher} & \textbf{Conferido?} \\
\midrule
\endhead
Entidade / ID do participante & & $\square$ \\
F1 --- Persona (P0\_) & & $\square$ \\
F2 --- Evento (trecho do nome) & & $\square$ \\
F2 --- Ingresso A (validar por QR): associado fictício, CPF, código \texttt{\#\_\_\_} & &
  $\square$ \\
F2 --- Ingresso B (validar por código digitado): associado fictício, CPF, código \texttt{\#\_\_\_}
  & & $\square$ \\
F3 --- Associado inadimplente reservado (ativo): nome, CPF & & $\square$ \\
\bottomrule
\end{longtable}
\end{center}

\textit{Como preparar:} emitir, no painel (ou no app, logado como associado fictício do seed), 1
ingresso por associado fictício no evento escolhido para os ingressos A e B; anotar CPF e código.
Deixar o app do moderador logado na conta do ingresso A, com o QR aberto (ou imprimir o QR). Para
o inadimplente, em \textbf{Mensalidades → Relatório de inadimplência}, escolher um associado
listado que não seja usado em outro kit; anotar nome e CPF.

\subsection*{Kit --- papel Associado}

\begin{center}
\small
\begin{longtable}{p{7.5cm}p{3.3cm}p{2cm}}
\toprule
\textbf{Item} & \textbf{Preencher} & \textbf{Conferido?} \\
\midrule
\endhead
Entidade / ID do participante & & $\square$ \\
F1 --- Variante (A auto-cadastro / B vincular conta) & & $\square$ \\
F1-A --- Persona (P0\_) & & $\square$ \\
F1-B --- CPF pré-cadastrado (ativo, sem login) & & $\square$ \\
F2/F3 --- Conta de teste (e-mail e senha) & & $\square$ \\
F2 --- Evento (trecho do nome) & & $\square$ \\
\bottomrule
\end{longtable}
\end{center}

\textit{Como preparar:} para F1-B, cadastrar no painel, antes do dia, uma pessoa fictícia com CPF
novo e categoria de sócio definida (anotar o CPF, sem criar conta de login). Para a conta de teste
(F2/F3), usar um associado ativo, com categoria de sócio definida e login já existente, garantindo
uma mensalidade em aberto no mês --- é a conta que o moderador usa no app depois do Fluxo 1 (uma
conta por participante, já que pagar a mensalidade a consome). Conferir que o evento do F2 está
publicado e que a conta de teste ainda não tem ingresso nele.

\section*{7. Condução por fluxo (uso exclusivo do moderador)}

Para cada tarefa: dizer a fala indicada, deixar o participante executar a folha do fluxo e rodar a
verificação em seguida (\texttt{cd backend} antes). Estado necessário = item do kit (Seção 6).

\subsection*{Fluxo 1 --- Cadastro de associado}

\textbf{Diretoria.} Estado: ficha de persona ainda não usada. Fala: ``Alguém novo quer se associar à
entidade. Você recebeu esta ficha preenchida à mão --- cadastre essa pessoa no sistema como se
fosse um caso real.''
\begin{verbatim}
node scripts/verificar-tarefas.js 1 <CPF-DA-FICHA>
\end{verbatim}
Esperado: associado existe e está \texttt{ativo}. A conta de login não aparece vinculada (normal:
só a diretoria cadastrou).

\textbf{Associado, situação A (auto-cadastro).} Estado: ficha de persona ainda não usada; app
deslogado. Fala: ``Você quer se tornar associado da entidade e está criando sua conta pelo
aplicativo pela primeira vez.''
\begin{verbatim}
node scripts/verificar-tarefas.js 1 <CPF-DA-FICHA>
\end{verbatim}
Esperado: associado com status \texttt{pendente\_validacao} e conta de login vinculada.

\textbf{Associado, situação B (vincular conta).} Estado: CPF pré-cadastrado (ativo, sem login) no
cartão do kit; app deslogado. Fala: ``A diretoria já cadastrou você como associado, mas você nunca
entrou no aplicativo. Use o CPF que consta neste cartão para criar seu acesso.''
\begin{verbatim}
node scripts/verificar-tarefas.js 1 <CPF-DO-CARTAO>
\end{verbatim}
Esperado: status \texttt{ativo} e conta de login vinculada.

Cada Associado faz \textbf{uma} situação. Depois do Fluxo 1, sair dessa conta e entrar na
\textbf{conta de teste do kit} antes de devolver o celular (o login não é tarefa avaliada). A conta
criada no Fluxo 1 não é usada nos Fluxos 2 e 3.

\subsection*{Fluxo 2 --- Ingressos e validação de entrada}

\textbf{Associado.} Estado: conta de teste já logada; evento publicado com ingresso avulso. O QR
dele não é lido por ninguém (a validação é a tarefa da Diretoria, com ingressos à parte).
\begin{small}
\begin{verbatim}
node scripts/verificar-tarefas.js 2 <CPF-DA-CONTA-DE-TESTE> \
  "<trecho do nome do evento>"
\end{verbatim}
\end{small}
Esperado: ingresso emitido; entrada ainda não validada.

\textbf{Diretoria.} Estado: 2 ingressos pré-emitidos, em nome de associados fictícios --- A (ler por
câmera; QR no celular do moderador ou impresso) e B (digitar o código, anotado no cartão). Fala:
``Estes dois convidados chegaram ao baile. Confira a entrada de cada um.''
\begin{small}
\begin{verbatim}
node scripts/verificar-tarefas.js 2 <CPF-ASSOCIADO-FICTICIO-A> \
  "<trecho do nome do evento>"
node scripts/verificar-tarefas.js 2 <CPF-ASSOCIADO-FICTICIO-B> \
  "<trecho do nome do evento>"
\end{verbatim}
\end{small}
Esperado: ingresso emitido e entrada validada, com horário preenchido, nos dois. Se a câmera falhar
(permissão, luz, rede), registrar o motivo e pedir o ingresso A pelo caminho manual, contando como
intervenção; o ingresso B continua reservado para o caminho manual.

\subsection*{Fluxo 3 --- Mensalidade}

\textbf{Associado.} Estado: conta de teste com mensalidade em aberto. A folha já avisa que o
pagamento online é simulado (RF19): a tarefa avalia a interface e o registro no sistema, não um
pagamento real.
\begin{verbatim}
node scripts/verificar-tarefas.js 3 <CPF-DA-CONTA-DE-TESTE>
\end{verbatim}
Esperado: mensalidade paga, com data e forma de pagamento preenchidas.

\textbf{Diretoria.} Estado: associado ativo e inadimplente reservado (nome e CPF no kit); informar o
nome ao participante. Fala: ``Este associado veio até a sede pagar a mensalidade em atraso.
Registre o pagamento.''
\begin{verbatim}
node scripts/verificar-tarefas.js 3 <CPF-DO-ASSOCIADO-INADIMPLENTE>
\end{verbatim}
Esperado: mensalidade paga, com \texttt{forma\_pagamento = presencial}.

\section*{8. Depois das tarefas: SUS e entrevista}

\subsection*{8.1 Questionário SUS}

Aplicar o formulário correspondente ao papel do participante, imediatamente após o bloco de
tarefas: Google Forms, preenchido pelo próprio participante (ou pelo moderador, lendo os itens).
Escala de 1 (discordo totalmente) a 5 (concordo totalmente), 10 itens. A pontuação (0--100, fórmula
de Brooke, 1996) é calculada na coluna O da planilha de respostas de cada formulário.

\begin{itemize}[leftmargin=6mm]
  \item \textbf{Painel (Diretoria):} formulário
    \url{https://docs.google.com/forms/d/e/1FAIpQLSdnIhOWGjHlwoPrQDRv5Wzp0uLf5O1_GWcuPxeGDMc1s3m-xw/viewform};
    planilha
    \url{https://docs.google.com/spreadsheets/d/1QZ78az8-krw79-We2ifDlJJ7GC2hwSr3nSY7NVuG5tA/edit}
  \item \textbf{Aplicativo (Associado):} formulário
    \url{https://docs.google.com/forms/d/e/1FAIpQLSeqf1V7mwzuHNQGOUnb9qju-9fiCPsCne48TkF6PFQgSi6t8Q/viewform};
    planilha
    \url{https://docs.google.com/spreadsheets/d/1cmmUhJ9u3hmNuvuSCGTmISsNt7vZI4erwR3HJkKJ1cA/edit}
\end{itemize}

\subsection*{8.2 Roteiro de entrevista semiestruturada}

Perguntas abertas, conduzidas após o SUS:

\begin{enumerate}[leftmargin=6mm]
  \item De forma geral, como foi a experiência de usar o sistema hoje?
  \item Em qual momento das tarefas você se sentiu mais em dúvida ou perdido? O que ajudaria ali?
  \item Alguma tela ou termo usado no sistema não ficou claro? O vocabulário (categoria de sócio,
    cargos da diretoria, tipos de evento) fez sentido para a entidade, ou soou genérico demais?
  \item Comparando com o que vocês já faziam antes (papel, planilha, WhatsApp), esse processo pelo
    sistema parece mais útil, mais fácil, ou não fez diferença?
  \item (Só para quem usou o fluxo mediado) Você confiaria em pedir para alguém da diretoria fazer
    isso por você, caso preferisse não usar o sistema diretamente?
  \item (Só para diretoria) As categorias de sócio e o croqui do salão, do jeito que ficaram
    configurados, correspondem a como a entidade de fato organiza isso hoje?
  \item (Só para diretoria) Nos fluxos em que a diretoria age em nome do associado (como o
    cadastro feito pelo painel), a forma como o sistema trata essa aprovação corresponde a como a
    entidade decide isso hoje, ou falta algum nível de autorização interna?
  \item Tem alguma coisa que você esperava encontrar e não encontrou, ou que faria diferente?
\end{enumerate}

\section*{9. Se uma tarefa falhar}

\begin{enumerate}[leftmargin=6mm]
  \item Aplicar a regra de intervenção (depois de \textasciitilde2 min travado, indicar o
    campo/botão específico e contar como 1 intervenção). Se mesmo assim não avançar, parar a
    tarefa.
  \item Registrar \textit{não concluída} e o motivo (erro do participante / bug do sistema /
    rede). Bug do sistema vai para a lista de achados, não é apagado.
  \item Seguir para a próxima tarefa \textbf{sem refazer a anterior}: a próxima já tem seu estado
    preparado no kit. Nunca ``consertar'' o estado da tarefa seguinte com dados da que falhou.
\end{enumerate}

\section*{10. Depois da sessão}

\begin{enumerate}[leftmargin=6mm]
  \item Rodar os comandos de verificação de cada fluxo executado (se ainda não rodados durante a
    sessão) e anotar o resultado na planilha de observação, cruzando com o número de intervenções
    do moderador --- essa combinação classifica a tarefa como \textit{concluída sem ajuda} /
    \textit{concluída com intervenção} / \textit{não concluída}.
  \item Aplicar o formulário SUS correspondente ao papel do participante (se ainda não aplicado).
  \item Conduzir a entrevista semiestruturada (Seção 8.2).
  \item Marcar o kit como \textbf{usado} (personas, CPF pré-cadastrado, ingressos, inadimplente) e,
    se houver participante seguinte, abrir o próximo kit --- nunca reaproveitar itens.
\end{enumerate}

\end{document}
